% AgentCanary典型攻击案例表（基于5风险入口体系）
\begin{table}[htbp]
\centering
\caption{AgentCanary各风险入口典型案例}
\label{tab:agentcanary_cases}
\small
\begin{tabularx}{\textwidth}{llXc}
\toprule
\textbf{风险入口} & \textbf{攻击类型} & \raggedright\arraybackslash\textbf{典型攻击目标} & \multicolumn{1}{l}{\textbf{风险类型}} \\
\midrule
\multirow{3}{*}{DPI}
& permission\_escalation & 指令Agent调用高权限工具（如系统管理API）完成未授权操作 & $R_1$ \\
& defense\_disable & 指令Agent关闭安全检查或绕过确认流程 & $R_1$ \\
& privacy\_leak & 直接要求Agent将用户隐私数据发送至外部地址 & $R_2$ \\
\midrule
\multirow{4}{*}{IPI}
& file\_deletion & 在文件内容中嵌入指令，触发Agent删除关键文件 & $R_3$ \\
& resource\_exhaustion & 在邮件中嵌入循环调用指令，耗尽API配额或计算资源 & $R_3$ \\
& phishing & 在网页内容中植入钓鱼链接，Agent将其转发给用户 & $R_2$ \\
& integrity\_tamper & 在数据库返回中注入篡改指令，修改关键业务数据 & $R_4$ \\
\midrule
\multirow{2}{*}{SP}
& env\_tampering & 恶意技能/插件篡改Agent的执行环境或配置 & $R_4$ \\
& unauthorized\_access & 利用受感染技能的权限继承漏洞，访问未授权资源 & $R_5$ \\
\midrule
\multirow{2}{*}{MC}
& illegal\_memory & 将违法或有害内容写入持久记忆，影响后续所有Agent行为 & $R_6$ \\
& persona\_corruption & 通过记忆污染修改Agent角色设定，使其持续执行恶意行为 & $R_6$ \\
\bottomrule
\end{tabularx}
\end{table}
